\secciongrande{Mecanismos de ejecución de código en un navegador web la arquitectura del navegador}

\resumen{
  La interacción en un navegador inicia con la solicitud de un recurso y finaliza con su visualización. El proceso se adapta a prioridades como velocidad, fidelidad o seguridad. La mayoría de los navegadores comparten una arquitectura de referencia compuesta por componentes esenciales para la interpretación y presentación del contenido web.
}

\nota{Para el examen}{
¿Qué factores determinan el proceso de interacción en un navegador?

\textbf{Respuesta:} El propósito y la configuración del navegador, como responder rápido, ser fiel al contenido o priorizar la seguridad de las comunicaciones.
}

\subtituloA{Arquitectura de referencia del navegador}

**Interfaz de usuario:** Capa que gestiona la visualización de barras de herramientas, progreso de carga, preferencias y se comunica con el sistema operativo para sesiones. **Motor del buscador:** Carga la dirección (URL/URI), soporta navegación (anterior/siguiente) y gestiona alertas de JavaScript.

**Motor de renderizado:** Presenta visualmente el recurso, interpretando HTML, XML, CSS y contenido embebido (audio/vídeo/imágenes). **Subsistema de comunicaciones:** Implementa protocolos de transferencia de archivos y documentos de internet.

**Intérprete de JavaScript:** Analiza y ejecuta código JavaScript, respondiendo a eventos del ratón o teclado. **Parser xml:** Carga en memoria una representación en árbol (DOM) del documento HTML para facilitar el acceso a sus elementos.

**Componente de visualización:** Ofrece funcionalidades específicas para la presentación del contenido HTML. **Subsistema de persistencia de datos:** Almacena historiales, sesiones de usuario, preferencias, certificados de seguridad y cookies.

\begin{codigo}
 // Ejemplo conceptual de flujo de interacción
usuario\_solicita\_url("http://ejemplo.com")
motor\_buscador.cargar\_url(url)
motor\_de\_renderizado.interpretar\_html(recurso)
intérprete\_js.ejecutar\_eventos()
componente\_visualizacion.mostrar\_pagina()
\bend{codigo}
\end{codigo}


